% ============================================================================
%  Problems: ความสัมพันธ์และฟังก์ชัน (Relations and Functions)
%  Ordered from easy (basic) → hard (POSN level).
%  Shared by exercise.tex and answer-key.tex
% ============================================================================

% --- Part 1: พื้นฐาน (Basic) ---
\examsection{ตอนที่ 1: พื้นฐาน}

\problem{
  กำหนด $A = \{1, 2\}$ และ $B = \{a, b\}$
  จงเขียนสมาชิกทั้งหมดของ $A \times B$
}{
  $A \times B = \{(1,a), (1,b), (2,a), (2,b)\}$
}

\problem{
  กำหนด $r = \{(1, 4), (2, 5), (3, 6), (4, 4)\}$
  จงหา $D_r$ (โดเมน) และ $R_r$ (เรนจ์) ของ $r$
}{
  $D_r = \{1, 2, 3, 4\}$,
  $R_r = \{4, 5, 6\}$
}

\problem{
  กำหนดความสัมพันธ์ $r = \{(x, y) \mid y = x^2,\; x \in \{-2, -1, 0, 1, 2\}\}$
  จงเขียนสมาชิกทั้งหมดของ $r$ พร้อมทั้งหา $D_r$ และ $R_r$
}{
  $r = \{(-2,4), (-1,1), (0,0), (1,1), (2,4)\}$,
  $D_r = \{-2, -1, 0, 1, 2\}$,
  $R_r = \{0, 1, 4\}$
}

\problem{
  กำหนด $f(x) = 3x - 5$ จงหาค่าของ
  \begin{enumerate}
    \item[(ก)] $f(2)$
    \item[(ข)] $f(0)$
    \item[(ค)] $f(-1)$
    \item[(ง)] $f(a)$ เมื่อ $a$ เป็นจำนวนจริงใด ๆ
  \end{enumerate}
}{
  (ก) $f(2) = 1$,
  (ข) $f(0) = -5$,
  (ค) $f(-1) = -8$,
  (ง) $f(a) = 3a - 5$
}

\problem{
  กำหนด $f(x) = \dfrac{x - 1}{3}$ จงหา $f^{-1}(x)$
  (ฟังก์ชันผกผันของ $f$)
}{
  $f^{-1}(x) = 3x + 1$
}

% --- Part 2: การประยุกต์ (Intermediate) ---
\examsection{ตอนที่ 2: การประยุกต์}

\problem{
  จงตรวจสอบว่า $r = \{(1, 2), (2, 3), (1, 4), (3, 5)\}$
  เป็นฟังก์ชันหรือไม่ เพราะเหตุใด
}{
  ไม่เป็นฟังก์ชัน เพราะ $1$ จับคู่กับทั้ง $2$ และ $4$ (สมาชิกตัวเดียวกันในโดเมนจับคู่กับสมาชิกในเรนจ์มากกว่า 1 ตัว)
}

\problem{
  จงหาโดเมนของฟังก์ชันต่อไปนี้
  \begin{enumerate}
    \item[(ก)] $f(x) = \dfrac{1}{x - 3}$
    \item[(ข)] $g(x) = \sqrt{x + 2}$
  \end{enumerate}
}{
  (ก) $D_f = \R - \{3\}$,
  (ข) $D_g = \{x \in \R \mid x \geq -2\} = [-2, \infty)$
}

\problem{
  กำหนด $f(x) = x^2 + 1$ จงหาเรนจ์ของ $f$
  (Hint: $x^2 \geq 0$ สำหรับทุกจำนวนจริง $x$)
}{
  $R_f = \{y \in \R \mid y \geq 1\} = [1, \infty)$
}

\problem{
  กำหนด $f(x) = 2x + 1$ และ $g(x) = x^2$
  จงหา
  \begin{enumerate}
    \item[(ก)] $(f \circ g)(3)$
    \item[(ข)] $(g \circ f)(x)$
    \item[(ค)] $(f \circ g)(x)$
  \end{enumerate}
}{
  (ก) $(f \circ g)(3) = f(9) = 19$,
  (ข) $(g \circ f)(x) = (2x + 1)^2 = 4x^2 + 4x + 1$,
  (ค) $(f \circ g)(x) = 2x^2 + 1$
}

\problem{
  กำหนด $f(2-x) = 2x + 1$ จงหา
  \begin{enumerate}
    \item[(ก)] $f(2)$
    \item[(ข)] $f(3)$
    \item[(ค)] $f(x)$
  \end{enumerate}
}{
  (ก) $f(2)=1$,
  (ข) $f(3)=-1$,
  (ค) $f(x)=5-2x$
}

% --- Part 3: ระดับ สอวน. (POSN Level) ---
\examsection{ตอนที่ 3: ระดับข้อสอบ สอวน.}

\problem{
  กำหนด $f: \R \to \R$ นิยามโดย $f(x) = 2x - 3$
  จงตรวจสอบว่า $f$ เป็นฟังก์ชัน \textbf{หนึ่งต่อหนึ่ง (1-1)} หรือไม่
  พร้อมทั้งตรวจสอบว่าเป็นฟังก์ชัน \textbf{ทั่วถึง (Onto)} หรือไม่
  (Hint: แก้สมการ $f(a_1) = f(a_2)$ เพื่อทดสอบ 1-1
  และแก้สมการ $f(x) = b$ สำหรับ $b \in \R$ ใด ๆ เพื่อทดสอบ Onto)
}{
  $f$ เป็นฟังก์ชัน 1-1 เพราะ $2a_1 - 3 = 2a_2 - 3 \implies a_1 = a_2$
  $f$ เป็นฟังก์ชัน Onto เพราะสำหรับทุก $b \in \R$ มี $x = \frac{b+3}{2} \in \R$ ที่ทำให้ $f(x) = b$
  ดังนั้น $f$ เป็นฟังก์ชัน Bijective
}

\problem{
  กำหนด $A = \{1, 2, 3\}$ และความสัมพันธ์ $r = \{(1,1), (2,2), (3,3), (1,2), (2,1)\}$ บน $A$
  จงตรวจสอบว่า $r$ มีสมบัติต่อไปนี้หรือไม่
  \begin{enumerate}
    \item[(ก)] สะท้อน (Reflexive)
    \item[(ข)] สมมาตร (Symmetric)
    \item[(ค)] ถ่ายทอด (Transitive)
  \end{enumerate}
  และ $r$ เป็นความสัมพันธ์สมมูล (Equivalence Relation) หรือไม่
}{
  (ก) Reflexive $\checkmark$ — มี $(1,1),(2,2),(3,3)$ ครบทุกตัว

  (ข) Symmetric $\checkmark$ — $(1,2)$ มี $(2,1)$, $(2,1)$ มี $(1,2)$

  (ค) Transitive $\checkmark$ — ตรวจสอบทุกกรณีแล้วผ่านทั้งหมด

  ดังนั้น $r$ เป็น Equivalence Relation $\checkmark$
}

\problem{
  (ข้อสอบจริงปี 68 ข้อ 8) กำหนด $f(x)=2(ax)^2-ax-1$ และ $g(x)=ax-1$ โดยที่ $a>0$ และ $x \neq \frac{1}{a}$ ถ้า $\frac{f}{g}(2a)=2$ แล้วจงหาค่าของ $f \circ g^{-1}(a)$
}{
  $f \circ g^{-1}(a) = 2$
}

\problem{
  (ข้อสอบจริงปี 65 ข้อ 39) จงพิจารณาข้อความต่อไปนี้ 
  \begin{enumerate}
    \item[1)] $f(x)=-3x^2-6x-8$ เป็นฟังก์ชัน $1-1$
    \item[2)] $f(x)=\{(2,-1),(-5,-5),(6,-1)\}$  เป็นฟังก์ชันจาก $B$ ไปทั่วถึง $A$ เมื่อ $A=\{-1,2,-5\}, B=\{2,-5,6\}$ 
  \end{enumerate}
  ข้อใดถูกต้องบ้าง
}{
  ช้อ 1 และ ข้อ 2 ผิด
}

\problem{
  (ข้อสอบจริงปี 65 ข้อ 40) กำหนด $f(x)=3x+5$ และ $h(x)=3x^2+3x-1$ \newline
  ถ้า $g$ เป็นฟังก์ชัน ซึ่งทำให้ $f \circ g=h$ แล้ว $g(5)$ มีค่าเท่าใด
}{
  $g(5)=28$
}

